% !TeX program = pdflatex
%==============================================================================
%  Avaliação Teórica 01 --- Aulas 01 a 06 (Bloco I)
%
%  O gabarito sai de "Avaliacao teorica 01 - gabarito.tex", que apenas liga a
%  opção [gabarito] e inclui este arquivo --- fonte única, dois PDFs.
%
%  AS TRÊS QUESTÕES FORAM SORTEADAS da "Lista de revisao 01-06", com cobertura
%  (uma aula por questão). O sorteio é reproduzível:
%
%      import numpy as np
%      rng = np.random.default_rng(20260914)      # a data, aaaammdd
%      aulas = sorted(rng.choice([1,2,3,4,5,6], size=3, replace=False).tolist())
%      # -> [2, 3, 6]; depois um exercício dentro de cada aula, na mesma ordem
%
%  Resultado:
%      Aula 02 -> Exercício  5  "a Ridge como moda a posteriori"
%      Aula 03 -> Exercício  7  "a mesma validação cruzada, escrita de dois jeitos"
%      Aula 06 -> Exercício 13  "para quem a régua importa"
%
%  NÃO HÁ TEXTO DE EXERCÍCIO NESTE ARQUIVO. Os enunciados e as soluções vêm por
%  \input dos MESMOS arquivos que a lista de revisão usa, em exercicios/ --- é
%  isso que garante que os dois documentos nunca divirjam. Aqui ficam só o
%  cabeçalho da prova, os títulos e pesos das questões, e os critérios de
%  correção, que são específicos da avaliação.
%
%  O \provatrue abaixo liga a pontuação por item: o comando \pts{1,0}, escrito
%  dentro dos arquivos de exercício, some na lista e aparece aqui.
%==============================================================================
\providecommand{\gabaritoopt}{}
\documentclass[11pt,a4paper]{article}
\usepackage[\gabaritoopt]{../recursos/latex/estilo-avaliacao}

\provatrue

\begin{document}
\cabecalhoprova{01}{01 a 06}{Fundamentos, Regressão, Validação Cruzada, KNN,
Árvores e \emph{Pipelines}}{2 horas}{12,0 pontos}

%==============================================================================
\questao{a Ridge como moda a posteriori}{4,0 pontos}
\daaula{02}
%==============================================================================
% Aula 02 --- a Ridge como moda a posteriori sob priori gaussiana.
Já mencionamos que o Ridge tem uma leitura bayesiana: ele é a moda a
posteriori sob uma priori gaussiana. A afirmação é mais do que curiosidade ---
ela dá um significado concreto ao $\lambda$, que até aqui era só um hiperparâmetro.

Considere o modelo
\[
  \bm y \mid \bbeta \;\sim\; N\big(\mathbb{X}\bbeta,\ \sigma^2\mathbb{I}\big),
  \qquad
  \bbeta \;\sim\; N\big(\bm 0,\ \sigma_\beta^2\mathbb{I}\big),
\]
com $\sigma^2$ e $\sigma_\beta^2$ conhecidos.
\begin{itens}
  \item \pts{1,6} Escreva a densidade a posteriori de $\bbeta$ a menos de
        constante multiplicativa e mostre que maximizá-la equivale a minimizar
        \[
          \frac{1}{2\sigma^2}\norm{\bm y-\mathbb{X}\bbeta}^2
          + \frac{1}{2\sigma_\beta^2}\norm{\bbeta}^2 .
        \]
        O caminho é o de sempre: passar ao logaritmo transforma produto em soma, e
        trocar o sinal transforma maximizar em minimizar.
  \item \pts{1,0} Compare o que você obteve com a definição do Ridge da aula e
        identifique $\lambda$ em função de $\sigma^2$ e $\sigma_\beta^2$. Cuidado
        com o fator que precisa ser eliminado antes da comparação.
  \item \pts{1,4} Agora a parte interessante: o que a priori está \emph{afirmando}
        quando $\sigma_\beta^2\to\infty$, e quando $\sigma_\beta^2\to 0$? Ligue
        cada resposta ao comportamento do Ridge em $\lambda\to 0$ e
        $\lambda\to\infty$.
\end{itens}

\begin{solucao}
\textbf{(a)} Pelo teorema de Bayes, a posteriori é proporcional ao produto da
verossimilhança pela priori:
\[
  \pi(\bbeta\mid\bm y)\;\propto\; f(\bm y\mid\bbeta)\,\pi(\bbeta)
   \;\propto\;
   \exp\Big\{-\frac{1}{2\sigma^2}\norm{\bm y-\mathbb{X}\bbeta}^2\Big\}
   \cdot
   \exp\Big\{-\frac{1}{2\sigma_\beta^2}\norm{\bbeta}^2\Big\}.
\]
As constantes de normalização das duas gaussianas não dependem de $\bbeta$ e por
isso podem ser absorvidas no $\propto$.

Maximizar uma função positiva é o mesmo que maximizar seu logaritmo, e o
logaritmo de um produto de exponenciais é a soma dos expoentes:
\[
  \log\pi(\bbeta\mid\bm y)
   = \text{const} - \frac{1}{2\sigma^2}\norm{\bm y-\mathbb{X}\bbeta}^2
     - \frac{1}{2\sigma_\beta^2}\norm{\bbeta}^2 .
\]
Maximizar isso é minimizar o simétrico dos dois termos que dependem de $\bbeta$,
que é exatamente o objetivo do enunciado.

\smallskip
\textbf{(b)} O objetivo de (a) tem um $1/(2\sigma^2)$ na frente do primeiro termo,
e a definição do Ridge não tem fator nenhum ali. Multiplique tudo por
$2\sigma^2>0$ --- o que não altera o minimizador, por ser constante positiva:
\[
  \norm{\bm y-\mathbb{X}\bbeta}^2 + \frac{\sigma^2}{\sigma_\beta^2}\norm{\bbeta}^2 .
\]
Agora a comparação é imediata, e
\[
  \boxed{\;\lambda=\frac{\sigma^2}{\sigma_\beta^2}\;}
\]

\smallskip
\textbf{(c)} Quando $\sigma_\beta^2\to\infty$, a priori fica \textbf{difusa}: ela
espalha massa por toda parte e não expressa preferência alguma por coeficientes
pequenos. Em outras palavras, ela não afirma nada. Correspondentemente
$\lambda\to 0$ e o Ridge degenera no MQO: os dados decidem sozinhos, com viés
baixo e variância alta.

Quando $\sigma_\beta^2\to 0$, a priori fica \textbf{concentrada em zero}: ela
afirma, antes de ver dado nenhum, que os coeficientes são praticamente nulos.
Correspondentemente $\lambda\to\infty$ e $\widehat\bbeta\to\bm 0$: viés máximo,
variância mínima.

Entre os dois extremos, a fórmula $\lambda=\sigma^2/\sigma_\beta^2$ diz uma coisa
bem razoável: regulariza-se mais quando os \emph{dados} são ruidosos
($\sigma^2$ grande) ou quando a \emph{crença prévia} é firme ($\sigma_\beta^2$
pequeno). O $\lambda$ é a razão entre duas incertezas, e escolhê-lo por validação
cruzada é, nessa leitura, deixar os dados escolherem a priori que melhor prediz.
\end{solucao}


\begin{criterio}
\textbf{(a)} 0,6 pelo produto verossimilhança $\times$ priori; 0,6 por passar ao
logaritmo e trocar máximo por mínimo corretamente. Ignorar as constantes de
normalização não desconta, desde que o aluno justifique que não dependem de
$\bbeta$.

\textbf{(b)} 0,4 por eliminar o fator $1/(2\sigma^2)$ sem alterar o argmin; 0,4
pelo valor de $\lambda$. Chegar a $\lambda$ invertido
($\sigma_\beta^2/\sigma^2$) vale 0,2.

\textbf{(c)} 0,4 por extremo. Exige ligar priori e $\lambda$, não apenas
descrever o comportamento do Ridge. Mencionar o par viés--variância é esperado,
mas não obrigatório para a nota cheia.

\textbf{(d)} 0,6 por identificar que a esparsidade é propriedade da moda; 0,6 por
explicar que a posteriori é contínua e a média integra. Aceitar qualquer
formulação correta do argumento de continuidade.
\end{criterio}

%==============================================================================
\questao{a mesma validação cruzada, escrita de dois jeitos}{4,0 pontos}
\daaula{03}
%==============================================================================
% Aula 03 --- as duas formas da fórmula de k dobras (aula x ISLP).
Quem estuda validação cruzada pelas notas de aula e pelo livro encontra duas fórmulas
diferentes para a mesma quantia, e é natural desconfiar de que uma delas esteja
errada. Não está: elas só diferem na unidade que cada uma pondera igualmente.

As notas de aula definem
\[
  \riscohat_{k\text{-CV}}
   = \frac1n\sum_{j=1}^{k}\sum_{i\in L_j}\big(Y_i-g_{-j}(\X_i)\big)^2 ,
\]
e o livro define $\mathrm{CV}_{(k)}=\frac1k\sum_{j=1}^{k}\mathrm{EQM}_j$, com
$\mathrm{EQM}_j=\frac{1}{\abs{L_j}}\sum_{i\in L_j}(Y_i-g_{-j}(\X_i))^2$.
\begin{itens}
  \item \pts{1,2} Reescreva a fórmula das notas de modo a fazer os
        $\mathrm{EQM}_j$ aparecerem, mostrando que
        \[
          \riscohat_{k\text{-CV}}=\sum_{j=1}^{k}\frac{\abs{L_j}}{n}\,\mathrm{EQM}_j .
        \]
  \item \pts{0,8} Olhando para os pesos que você obteve, diga em que condição as
        duas definições coincidem --- e se coincidem exatamente ou apenas
        aproximadamente. Em seguida, nomeie a unidade que cada fórmula pondera
        igualmente.
  \item \pts{0,7} Com $n=50$ e $k=7$, as dobras não podem ter todas o mesmo
        tamanho; suponha que tenham tamanhos $8,8,7,7,7,7,6$. Confirme que somam $50$
        e explique por que, neste caso, os dois números diferem.
  \item \pts{1,3} Suponha que a dobra de $6$ observações tenha $\mathrm{EQM}=10$ e
        todas as outras tenham $\mathrm{EQM}=1$. Calcule os dois estimadores, diga
        qual deles dá mais peso à dobra atípica, e qual você reportaria. Justifique
        a escolha.
\end{itens}

\begin{solucao}
\textbf{(a)} A chave é reconhecer a soma interna. Pela definição de $\mathrm{EQM}_j$,
basta multiplicar dos dois lados por $\abs{L_j}$:
\[
  \sum_{i\in L_j}\big(Y_i-g_{-j}(\X_i)\big)^2=\abs{L_j}\,\mathrm{EQM}_j .
\]
Substituindo na fórmula das notas,
\[
  \riscohat_{k\text{-CV}}
   =\frac1n\sum_{j=1}^{k}\abs{L_j}\,\mathrm{EQM}_j
   =\sum_{j=1}^{k}\frac{\abs{L_j}}{n}\,\mathrm{EQM}_j .
\]
São os \emph{mesmos} $\mathrm{EQM}_j$ do livro, combinados com pesos
$w_j=\abs{L_j}/n$. E esses pesos somam $1$, porque as dobras particionam as $n$
observações: $\sum_j\abs{L_j}=n$. As duas fórmulas são, portanto, médias ponderadas dos
mesmos $k$ números, e só diferem nos pesos: $\abs{L_j}/n$ numa, $1/k$ na outra.

\smallskip
\textbf{(b)} Os pesos coincidem em todas as dobras se e só se $\abs{L_j}/n=1/k$ para todo
$j$, isto é, se todas as dobras têm tamanho $n/k$ --- o que só é possível quando $k$
divide $n$. Nesse caso as duas fórmulas são a mesma soma dos mesmos números com os
mesmos pesos: coincidem \textbf{exatamente}, e não aproximadamente.

A unidade que cada uma pondera igualmente: as notas dão o mesmo peso a cada
\textbf{observação} --- todo resíduo de validação entra com peso $1/n$ ---; o livro dá o
mesmo peso a cada \textbf{dobra}. Visto do lado das observações, o livro dá a cada
resíduo da dobra $j$ o peso $1/(k\abs{L_j})$.

\smallskip
\textbf{(c)} $8+8+7+7+7+7+6=50$.

Os pesos das notas são $8/50$ (duas vezes), $7/50$ (quatro vezes) e $6/50$; os do
livro, $1/7$ para cada dobra. Como os tamanhos diferem, os pesos diferem, e as duas
médias só coincidem por acaso --- quando os $\mathrm{EQM}_j$ são todos iguais, por
exemplo.

Olhando por observação fica mais claro. No livro, cada resíduo recebe
\[
  \frac{1}{7\cdot8}=\frac1{56},\qquad
  \frac{1}{7\cdot7}=\frac1{49}\qquad\text{ou}\qquad
  \frac{1}{7\cdot6}=\frac1{42},
\]
conforme a dobra em que caiu, contra $1/50$ para todos nas notas. Uma observação da
dobra de $6$ pesa $4/3$ de uma da dobra de $8$, só porque divide a sua dobra com menos
gente.

\smallskip
\textbf{(d)} Pelas notas, a dobra atípica entra com peso $6/50$ e as outras $44$
observações com $44/50$:
\[
  \riscohat_{k\text{-CV}}=\frac{6}{50}\times10+\frac{44}{50}\times1=\frac{104}{50}=\mathbf{2{,}08}.
\]
Pelo livro, a dobra atípica é uma das sete, e só:
\[
  \mathrm{CV}_{(7)}=\frac17\big(10+6\times1\big)=\frac{16}{7}\approx\mathbf{2{,}286}.
\]
O livro dá mais peso à dobra atípica: $1/7\approx0{,}143$ contra $6/50=0{,}12$. Ela é a
menor dobra e, ainda assim, vale um sétimo do total.

\emph{Qual reportar?} Para \emph{estimar o risco}, a das notas tem a seu favor um
argumento limpo. O risco é a esperança da perda numa observação nova, e cada resíduo de
validação $\ell_i$ é uma medida dessa perda. Se essas medidas tivessem a mesma média e a
mesma variância $s^2$, e fossem não correlacionadas --- uma idealização, mas é ela que
orienta ---, toda combinação $\sum_iw_i\ell_i$ com pesos somando $1$ teria a mesma
média, e variância $s^2\sum_iw_i^2$. Pela desigualdade de Cauchy--Schwarz,
\[
  1=\Big(\sum_{i=1}^n w_i\Big)^2\le n\sum_{i=1}^n w_i^2 ,
\]
com igualdade só nos pesos iguais, $w_i=1/n$. Pesar as observações igualmente é o que
menos desperdiça informação; dar mais peso a quem caiu numa dobra pequena é deixar o
acaso do corte decidir.

A do livro também se defende. Ela trata cada dobra como uma repetição do experimento
"treinar e validar", e é a escala natural quando o que se quer é comparar as dobras
entre si: é dos $k$ números $\mathrm{EQM}_j$ que sai o erro-padrão da regra de um
erro-padrão. E é a que o \texttt{scikit-learn} calcula ---
\texttt{cross\_val\_score(\dots).mean()} é média não ponderada ---, então o essencial é
\emph{saber qual das duas o seu código calcula} antes de comparar o seu número com o de
outra pessoa.

E a recomendação das notas encerra a discussão pela raiz: com $n$ pequeno, escolha um
$k$ que divida $n$. Aí as duas coincidem, e a pergunta não se coloca.
\end{solucao}


\begin{criterio}
\textbf{(a)} 0,6 por reconhecer que
$\sum_{i\in L_j}(\cdot)^2=\abs{L_j}\mathrm{EQM}_j$; 0,6 pela substituição e pela
forma final. Observar que os pesos somam $1$ é bônus.

\textbf{(b)} 0,5 pela igualdade dos pesos quando $k$ divide $n$; 0,3 por nomear as
duas unidades (observação $\times$ dobra). Responder "aproximadamente iguais"
perde 0,2 --- a igualdade é exata.

\textbf{(c)} 0,2 pela soma; 0,5 pela explicação dos pesos desiguais. Basta o
argumento; não é preciso calcular os sete pesos.

\textbf{(d)} 0,45 por cada valor correto ($2{,}08$ e $\approx 2{,}286$); 0,4 pela
justificativa da escolha. Aceitar tanto a defesa da fórmula ponderada quanto a da
não ponderada, desde que argumentada; \textbf{não} aceitar resposta que ignore a
existência das duas.
\end{criterio}

%==============================================================================
\questao{para quem a régua importa}{4,0 pontos}
\daaula{06}
%==============================================================================
% Aula 06 --- quem é sensível à escala das covariáveis, e por quê.
Um colega recebeu um banco em que a coluna \texttt{renda} está em reais e a coluna
\texttt{altura} em metros, e pergunta se precisa padronizar antes de ajustar. A
resposta honesta é "depende do método" --- e o critério que decide cabe numa
pergunta: \textbf{o método soma coisas vindas de colunas diferentes?}

Se soma, as unidades precisam ser comparáveis, e trocar uma delas estraga a
comparação. Se não soma, a unidade é irrelevante.
\begin{itens}
  \item \pts{2,4 --- 0,4 cada} Para cada método abaixo, diga se multiplicar uma
        única coluna por $1000$ \textbf{muda} ou \textbf{não muda} a predição,
        justificando em uma linha onde está --- ou onde não está --- a soma sobre
        colunas.
        \begin{enumerate}[label=(\roman*),leftmargin=1.0cm,nosep]
          \item Mínimos quadrados ordinários.
          \item Regressão Ridge com $\lambda>0$ fixo.
          \item Árvore de regressão.
          \item Floresta aleatória.
          \item KNN com $k$ fixo.
          \item \emph{Spline} de base truncada, ajustado por mínimos quadrados.
        \end{enumerate}
  \item \pts{0,7} Entre os que mudam, qual você espera que seja o mais afetado, e
        por quê? Compare o \emph{modo} como a escala entra em cada um.
  \item \pts{0,9} Um colega observa que, ao reescalar uma coluna, o risco do Lasso
        \emph{melhorou}, e conclui que não padronizar foi uma boa ideia. Explique
        o erro do raciocínio.
\end{itens}

\begin{solucao}
\textbf{(a)}

\begin{center}\small
\renewcommand{\arraystretch}{1.25}
\begin{tabular}{@{}clp{8.2cm}@{}}
\toprule
 & muda? & onde está (ou não) a soma sobre colunas \\
\midrule
(i) & \textbf{não} & O MQO é equivariante: o coeficiente da coluna se divide por
$1000$ e o produto $\beta_jx_j$ fica igual. Nada na função-objetivo compara
coeficientes entre si. \\
(ii) & \textbf{sim} & A penalidade $\lambda\sum_j\beta_j^2$ \emph{soma}
coeficientes de colunas diferentes; mudar a unidade de uma delas muda quanto ela
pesa nessa soma. \\
(iii) & \textbf{não} & A árvore só pergunta "$x_j\le t$?". O corte $t$ é
multiplicado junto, e a pergunta separa exatamente as mesmas observações. \\
(iv) & \textbf{não} & É uma média de árvores, e cada uma é indiferente; a
indiferença atravessa a agregação. \\
(v) & \textbf{sim} & A distância euclidiana $\sum_j(x_{ij}-x_{\ell j})^2$ é uma
soma sobre colunas. \\
(vi) & \textbf{não} & É um MQO sobre uma base fixa de funções de \emph{uma}
coluna; vale o mesmo argumento de (i). \\
\bottomrule
\end{tabular}
\end{center}

\smallskip
\textbf{(b)} O \textbf{KNN}, e a diferença é de ordem, não de grau.

Nos métodos penalizados a escala entra \emph{uma vez}, de forma linear, na
comparação entre coeficientes. Na distância euclidiana ela entra \textbf{ao
quadrado}: multiplicar a coluna por $1000$ multiplica as diferenças naquela
coordenada por $1000$, e a contribuição delas para a distância por $10^6$.

O efeito é que a coluna reescalada passa a decidir sozinha quem é vizinho de quem.
As demais não são apenas menos importantes --- elas somem da conta, e o método
deixa de usar a informação que carregavam.

\smallskip
\textbf{(c)} O erro é confundir um resultado com um argumento.

A melhora aconteceu porque a coluna reescalada era, naquele problema, a de maior
efeito verdadeiro. Deixá-la $1000$ vezes maior tornou o coeficiente dela $1000$
vezes menor, e um coeficiente minúsculo quase não pesa em $\sum_j\abs{\beta_j}$ ---
a coluna ficou praticamente isenta da penalidade. Isentar de penalidade justamente
a coluna que mais importa calhou de ajudar.

Calhou: é sorte, não método. Aplique a mesma reescala a uma coluna
\emph{irrelevante} e o efeito se inverte --- ela ganharia liberdade para receber um
coeficiente grande sem pagar por isso, e o ajuste pioraria.

O ponto de fundo é que, nos dois casos, quem decidiu quanto regularizar cada
coluna não foi o analista: foi a \textbf{unidade de medida} em que o dado chegou.
E é por isso que o argumento do colega não se salva nem quando o resultado o
favorece --- um procedimento cujo desempenho depende de a coluna ter sido anotada
em metros ou em milímetros está fora de controle, e ter dado certo desta vez não
muda isso.
\end{solucao}


\begin{criterio}
\textbf{(a)} 0,2 pelo veredito e 0,2 pela justificativa, em cada um dos seis
itens. Veredito certo com justificativa ausente ou errada vale 0,2. Aceitar
"equivariante" sem demonstração em (i) e (vi).

\textbf{(b)} 0,3 por apontar o KNN; 0,4 pelo argumento do quadrado na distância.
Responder "Ridge" com bom argumento vale 0,3.

\textbf{(c)} 0,4 por identificar que a coluna reescalada era a de maior efeito;
0,5 por explicar que o efeito se inverteria numa coluna irrelevante. A observação
de que foi a unidade de medida que tomou a decisão é o ponto alto da resposta e
deve ser premiada dentro dos 0,5.
\end{criterio}

\vfill
\begin{center}\small\itshape Boa prova!\end{center}

\end{document}
